\documentclass{jsarticle}
\usepackage{ascmac,amsmath,amssymb,amsthm,bm,physics,arydshln,mathcomp}
\newtheorem{question}{問題}
\newtheorem*{answer*}{解答}
\begin{document}

\begin{question}
  $\displaystyle \int_\alpha^\beta (x-\alpha)(x-\beta) dx$ を求めよ.
\end{question}

\begin{question}
$\displaystyle \int_0^{\infty} \frac{x^2+1}{x^4+1} dx$ を求めよ.
\end{question}

\begin{question}
 $\displaystyle \int_0^1 \sqrt{\frac{1-x}{x}}dx$ を求めよ.
\end{question}

\begin{question}
$\displaystyle \int_0^{\frac{\pi}{2}} \frac{\sqrt{\sin x}}{\sqrt{\sin x} + \sqrt{\cos x}} dx$ を求めよ.
\end{question}

\begin{question}
$\displaystyle \int_0^{\infty} e^{-tx}\cdot \frac{\sin x}{x} dx \ (t>0)$を求めよ.
\end{question}

\begin{question}
$\displaystyle \int_0^{\pi} \log(1+\cos x) dx$ を求めよ.
\end{question}

\begin{question}
$\displaystyle \int_0^1 \frac{\arcsin x}{x} dx$ を求めよ.
\end{question}

\begin{question}
不定積分$\displaystyle \int \frac{1}{\sin^3 x} dx$ を求めよ.
\end{question}

\begin{question}
不定積分$\displaystyle \int \frac{1+x^2}{(1-x^2)\sqrt{x^2+x+1}} dx \ \ (-1 < x < 1)$ を求めよ.
\end{question}

\bigskip
\bigskip
\bigskip
\bigskip
\bigskip
\bigskip
\bigskip
\bigskip
\bigskip
\bigskip
\bigskip
\bigskip
\bigskip
\bigskip
\bigskip
\bigskip
\bigskip
\bigskip
\bigskip
\bigskip
\bigskip
\bigskip
\bigskip
\bigskip
\bigskip
\bigskip
\bigskip
\bigskip
\bigskip
\bigskip
\bigskip
\bigskip


\begin{answer*}
\end{answer*}


\begin{align*}
(1) (積分は工夫して解け) \displaystyle \int_\alpha^\beta (x-\alpha)(x-\beta) dx
&= \displaystyle \int_\alpha^\beta \{(x-\alpha)^2 + (\alpha - \beta)(x - \alpha) \} dx \\
&= \left[\dfrac{1}{3}(x-\alpha)^3 + \dfrac{\alpha - \beta}{2} (x-\alpha)^2 \right]_\alpha^\beta \\
&= \dfrac{1}{3}(\beta - \alpha)^3 + \dfrac{\alpha-\beta}{2}(\beta - \alpha)^2 \\
&= -\dfrac{1}{6}(\beta - \alpha)^3
\end{align*}


\begin{align*}
(2) \hspace{4zw}
\displaystyle \int_0^{\infty} \frac{x^2+1}{x^4+1} dx
&= \displaystyle \frac{1}{2} \int_0^{\infty} \frac{(x^2 + \sqrt{2}x + 1) + (x^2 - \sqrt{2}x + 1)}{(x^2 + \sqrt{2}x + 1) (x^2 - \sqrt{2}x + 1)} dx\\
&= \displaystyle \frac{1}{2} \int_0^{\infty} \left(\frac{1}{x^2 - \sqrt{2}x + 1} + \frac{1}{x^2 + \sqrt{2}x + 1} \right) dx \\
&= \displaystyle \int_0^{\infty} \left\{\frac{1}{(\sqrt{2}x+1)^2 + 1} + \frac{1}{(\sqrt{2}x - 1)^2 + 1} \right\} \\
&= \displaystyle \frac{1}{\sqrt{2}} \left[\arctan(\sqrt{2}x + 1) + \arctan(\sqrt{2}x - 1) \right]_0^{\infty} \\
&= \displaystyle \frac{\pi}{\sqrt{2}} \\
&(\arctan(-1) = \frac{3}{4} \pi としないように注意)
\end{align*}

\bigskip
\bigskip
\bigskip
\bigskip
\bigskip

\begin{align*}
(3) \hspace{10zw} \displaystyle \int_0^1 \sqrt{\frac{1-x}{x}}dx 
&= \displaystyle \int_0^1 x^{-\frac{1}{2}} (1-x)^{\frac{1}{2}} dx \\
&= \displaystyle B\left(\frac{1}{2}, \ \frac{3}{2} \right) \\
&= \displaystyle \frac{\Gamma(\frac{3}{2}) \Gamma(\frac{1}{2})}{\Gamma(2)} \\
&= \displaystyle \frac{(\Gamma(\frac{1}{2}))^2}{2} \ \ (\because \Gamma(1+n) = n\Gamma(n)) \\
&= \displaystyle \frac{\pi}{2}
\end{align*}

\bigskip
\bigskip
\bigskip
\bigskip
\bigskip

\hspace{8zw} \textbf{別解}
\begin{align*}
  \displaystyle \int_0^1 \sqrt{\frac{1-x}{x}}dx
&= \displaystyle \int_0^1 (x)'\sqrt{\frac{1-x}{x}}dx \\
&= \left[x\sqrt{\frac{1-x}{x}} \right]_0^1 -  \displaystyle \int_0^1 x \left(\sqrt{\frac{1-x}{x}} \right)'dx \\
&= -\displaystyle \int_0^1 \frac{\left(\sqrt{\frac{1-x}{x}} \right)'dx}{1 + \left(\sqrt{\frac{1-x}{x}} \right)^2}dx \\
&= -\left[\arctan{\sqrt{\frac{1-x}{x}}} \right]_0^1 \\
&= \displaystyle \frac{\pi}{2}
\end{align*}

\begin{align*}
(4) \ t := \displaystyle \frac{\pi}{2} - x \ (King \ Property) と置換すると,
\displaystyle \int_0^{\frac{\pi}{2}} \frac{\sqrt{\sin x}}{\sqrt{\sin x} + \sqrt{\cos x}} dx 
&= \displaystyle \int_0^{\frac{\pi}{2}} \frac{\sqrt{\sin (\frac{\pi}{2} - t)}}{\sqrt{\sin (\frac{\pi}{2} - t)} + \sqrt{\cos (\frac{\pi}{2} - t)}} dt \\
&= \displaystyle \int_0^{\frac{\pi}{2}} \frac{\sqrt{\cos t}}{\sqrt{\cos t} + \sqrt{\sin t}} dt 
\end{align*}

\ \ \ なので, $2\displaystyle \int_0^{\frac{\pi}{2}} \frac{\sqrt{\sin x}}{\sqrt{\sin x} + \sqrt{\cos x}} dx
= \displaystyle \int_0^{\frac{\pi}{2}} \frac{\sqrt{\sin x} + \sqrt{\cos x}}{\sqrt{\sin x} + \sqrt{\cos x}} dx
= \displaystyle \frac{\pi}{2}$ \ となり

\ \ \ $\underline{\displaystyle \int_0^{\frac{\pi}{2}} \frac{\sqrt{\sin x}}{\sqrt{\sin x} + \sqrt{\cos x}} dx
= \displaystyle \frac{\pi}{4} }$

\bigskip
\bigskip
\bigskip
\bigskip
\bigskip
\bigskip
\bigskip

(5) $F(t) := \displaystyle \int_0^{\infty} e^{-tx}\cdot \frac{\sin x}{x} dx$ とおく.

\ \ \ \ \ $\displaystyle \left| \frac{\partial}{\partial t} \left(e^{-tx}\cdot \frac{\sin x}{x} \right) \right| = |e^{-tx} \sin x| \le e^{-tx}$ であり, $\displaystyle \int_0^{\infty} e^{-tx} dx = \frac{1}{t} < \infty$ なので, 
  
\ \ \ \ \ $\displaystyle \int_0^{\infty} \frac{\partial}{\partial t} \left(e^{-tx}\cdot \frac{\sin x}{x} \right) dx$ は一様収束する.
  
\begin{align*}
よって, F'(t)
&=\displaystyle \int_0^{\infty} \frac{\partial}{\partial t} \left(e^{-tx}\cdot \frac{\sin x}{x} \right) dx \\
&= -\displaystyle \int_0^{\infty} e^{-tx} \sin x dx \\
&= \displaystyle \frac{1}{t} \int_0^{\infty} (e^{-tx})' \sin x dx \\
&= \displaystyle \frac{1}{t} \left[e^{-tx} \sin x \right]_0^{\infty} - \displaystyle \frac{1}{t} \int_0^{\infty} e^{-tx} \cos x dx \\
&= \displaystyle \frac{1}{t^2} \int_0^{\infty} (e^{-tx})' \cos x dx \\
&= \displaystyle \frac{1}{t^2} \left[e^{-tx} \cos x \right]_{x=0}^{\infty} + \displaystyle \frac{1}{t^2} \int_0^{\infty} e^{-tx} \sin x dx \\
&= -\displaystyle \frac{1}{t^2} + \displaystyle \frac{1}{t^2} \int_0^{\infty} e^{-tx} \sin x dx \ \ \ \ \ なので
\end{align*}
\ \ \ \ \ $F'(t) = \displaystyle \frac{-1}{t^2+1}$. \ よって$F(t) = -\arctan t + C$ (Cは積分定数)

\ \ \ \ \ $\displaystyle \lim_{t \to \infty} F(t) = 0$ なので $C = \displaystyle \frac{\pi}{2}$. \\
\ \ \ \ \ したがって, $\underline{F(t) = \displaystyle \frac{\pi}{2} - \arctan t}$. \ \ \ \ \ $(t=0 を代入するとディリクレ積分となる)$

\bigskip
\bigskip
\bigskip
\bigskip
\bigskip
\bigskip
\bigskip

(6)
$t:= \pi - x$ と置換すると, 
$\displaystyle \int_0^{\pi} \log(1+\cos x) dx = \displaystyle \int_0^{\pi} \log(1-\cos t) dt$ なので

\begin{align*}
\displaystyle \int_0^{\pi} \log(1+\cos x) dx
&= \displaystyle \frac{1}{2} \int_0^{\pi} \{\log(1+\cos x) + \log(1-\cos x) \} dx \\
&= \displaystyle \int_0^{\pi} \log(\sin x) dx \\
&= \displaystyle \int_0^{\frac{\pi}{2}} \log(\sin x) dx + 
\displaystyle \int_{\frac{\pi}{2}}^{\pi} \log(\sin x) dx \\
&= -\displaystyle \frac{\pi}{2} \log2 + 
\displaystyle \int_{\frac{\pi}{2}}^{\pi} \log(\sin x) dx \\
&= -\displaystyle \frac{\pi}{2} \log2 + \displaystyle \int_0^{\frac{\pi}{2}} \log(\cos \theta) d\theta \ \ \ (\theta = x - \displaystyle \frac{\pi}{2} と置換) \\
&= \underline{-\pi\log2}
\end{align*}

$  \left(
  \begin{tabular}{l}
   ※ $\displaystyle \int_0^{\frac{\pi}{2}} \log\sin x dx$ については絶対収束することを示したうえで, $x = \displaystyle \frac{\pi}{2} - \theta, \ x = \pi - \theta$ と置換すれば求められる.
  \end{tabular}
  \right)$

\bigskip
\bigskip
\bigskip
\bigskip
\bigskip
\bigskip
\bigskip

\hspace{7zw} \textbf{別解}
\begin{align*}
\displaystyle \int_0^{\pi} \log(1+\cos x) dx
&= \displaystyle \int_0^{\pi} \log(2\cos^2\left(\frac{x}{2} \right)) dx \\
&= \pi\log2 + 2\displaystyle \int_0^{\pi} \log\cos{\frac{x}{2}} dx \\
&= \pi\log2 + 4\displaystyle \int_0^{\frac{\pi}{2}} \log\cos\theta \ d\theta \ \ \ \ (x = 2\theta と置換) \\
&= \underline{-\pi\log2}
\end{align*}

\bigskip
\bigskip
\bigskip
\bigskip
\bigskip
\bigskip
\bigskip

(7)
\begin{align*}
\displaystyle \int_0^1 \frac{\arcsin x}{x} dx
&= \displaystyle \int_0^1 \arcsin x (\log x)' dx \\
&= \left[\arcsin x \cdot \log x \right]_0^1 - \displaystyle \int_0^1 \frac{\log x}{\sqrt{1-x^2}} dx \\
&= -\displaystyle \int_0^{\frac{\pi}{2}} \log(\sin \theta) d\theta　\ \ \ \ (x = \sin \theta) \\
&= \underline{\displaystyle \frac{\pi}{2} \log2}
\end{align*}
\hspace{7zw}($\displaystyle \lim_{x \to 0} \arcsin x \cdot \log x
= \lim_{y \to 0} y\log{\sin y}
= \lim_{y \to 0} \frac{y}{\sin y} \sin y \cdot \log{\sin y}
= 1\cdot0 = 0$)

\bigskip
\bigskip
\bigskip
\bigskip
\bigskip
\bigskip
\bigskip

\hspace{7zw} \textbf{別解}
\begin{align*}
\displaystyle \int_0^1 \frac{\arcsin x}{x} dx
&= \displaystyle \int_0^{\frac{\pi}{2}} y\cdot \frac{\cos y}{\sin y} dy \\
&= \displaystyle \int_0^{\frac{\pi}{2}} y(\log{\sin y})' dy \\
&= \left[y\log{\sin y} \right]_0^{\frac{\pi}{2}} - \displaystyle \int_0^{\frac{\pi}{2}} \log\sin x dx \\
&= \underline{\displaystyle \frac{\pi}{2} \log2}
\end{align*}

\bigskip

\hspace{7zw} \textbf{別解２}
$\arcsin x$ をマクローリン展開.

\bigskip
\bigskip
\bigskip
\bigskip
\bigskip
\bigskip
\bigskip
\bigskip

(8)
$t := \displaystyle \tan{\frac{x}{2}}$ とおくと, 
$\displaystyle -\frac{\pi}{2} < \frac{x}{2} < \frac{\pi}{2}$ なので $\displaystyle \sin \frac{x}{2} = \frac{t}{\sqrt{1+t^2}}, \ \cos{\frac{x}{2}} = \frac{1}{\sqrt{1+t^2}}$ であり, 

$\sin x = \displaystyle \frac{2t}{1+t^2}, \ \cos x = \frac{1-t^2}{1+t^2}$.
\begin{align*}
\displaystyle \int \frac{1}{\sin^3 x} dx
&= \displaystyle \int \frac{(1+t^2)^3}{8t^3} \cdot \frac{2}{1+t^2} dt \\
&= \displaystyle \frac{1}{4} \int \left(t + \frac{2}{t} + t^{-3} \right) dt \\
&= \displaystyle \frac{1}{4} \left(\frac{1}{2} t^2 + 2\log|t| - \frac{1}{2}t^{-2} \right) + C \ \ (Cは積分定数) \\
&= \underline{\displaystyle \frac{1}{4} \left(\frac{1}{2} \displaystyle \tan^2{\frac{x}{2}} + 2\log \left|\displaystyle \tan{\frac{x}{2}} \right| - \frac{1}{2\displaystyle \tan^{2}{\frac{x}{2}}} \right) + C}
\end{align*}

\bigskip
\bigskip
\bigskip
\bigskip
\bigskip
\bigskip
\bigskip
\bigskip

(9)
\begin{align*}
\displaystyle \int \frac{1+x^2}{(1-x^2)\sqrt{x^2+x+1}} dx
&= \displaystyle \int \frac{-(1-x^2) + 2}{(1-x^2)\sqrt{x^2+x+1}} dx \\
&= -\displaystyle \int \frac{1}{\sqrt{x^2+x+1}} dx + \displaystyle \int \frac{2}{(1-x^2)\sqrt{x^2+x+1}} dx \\
&= -\displaystyle \int \frac{1}{\sqrt{x^2+x+1}} dx + \int \frac{1}{(1+x)\sqrt{x^2+x+1}} dx + \int \frac{1}{(1-x)\sqrt{x^2+x+1}} dx
\end{align*}

\begin{align*}
 \textcircled{\scriptsize 1} \displaystyle \int \frac{1}{\sqrt{x^2+x+1}} dx
&= \displaystyle \int \frac{1}{\sqrt{(x+\frac{1}{2})^2 + (\frac{\sqrt{3}}{2})^2}} dx \\
&= \log \left|x+\frac{1}{2} + \sqrt{x^2+x+1} \right| +
C_1 \ \ \ (C_1 は積分定数)
\end{align*}

\begin{align*}
 \textcircled{\scriptsize 2} \int \frac{1}{(1+x)\sqrt{x^2+x+1}} dx
 &= - \displaystyle \int \frac{1}{\sqrt{t^2 - t + 1}} dt \ \ (1+x = \frac{1}{t} と置換) \\
 &= - \displaystyle \int \frac{1}{\sqrt{(t - \frac{1}{2})^2 + (\frac{\sqrt{3}}{2})^2}} dt \\
 &= - \log \left| t-\frac{1}{2} + \sqrt{t^2-t+1}  \right| + C_2 \ \ \ (C_2 は積分定数) \\
 &= - \log \left| \frac{1-x+2\sqrt{x^2+x+1}}{2(x+1)} \right| + C_2
\end{align*}

\begin{align*}
\textcircled{\scriptsize 3} \int \frac{1}{(1-x)\sqrt{x^2+x+1}} dx
&= \displaystyle \int \frac{1}{\sqrt{3t^2-3t+1}} dt \ \ \left(1-x = \frac{1}{t}と置換 \right) \\
&= \displaystyle \frac{1}{\sqrt{3}} \int \frac{1}{\sqrt{(t-\frac{1}{2})^2 + \frac{1}{12}}} dt \\
&= \displaystyle \frac{1}{\sqrt{3}} \log \left| t-\frac{1}{2} + \sqrt{t^2 - t + \frac{1}{3}} \right| + C_3 \ \ \ (C_3 は積分定数)　\\
&= \displaystyle \frac{1}{\sqrt{3}} \log \left| \frac{\sqrt{3} (1 + x) + 2\sqrt{x^2 + x + 1}}{1-x} \right| + C_3
\end{align*}

したがって, \\
(与式)\\
= $\underline{-\displaystyle \log \left|x+\frac{1}{2} + \sqrt{x^2+x+1} \right| - \log \left| \frac{1-x+2\sqrt{x^2+x+1}}{2(x+1)} \right| + \displaystyle \frac{1}{\sqrt{3}} \log \left| \frac{\sqrt{3} (1 + x) + 2\sqrt{x^2 + x + 1}}{1-x} \right| + C}$

\hspace{41zw}($C$は積分定数)


\end{document}
